\documentclass[12pt]{article}
\usepackage[utf8]{inputenc}
\usepackage[T1]{fontenc}
\usepackage{hyperref}
\title{Takestani Language Resources for K-12 Curriculum Development}
\author{Compiled by Claude Sonnet 5 (claude-sonnet-5) for Andr\'as Kornai}
\date{}
\begin{document}
\maketitle

\section{Introduction}
Takestani (ISO 639-3 tks) is a variety of Southern Tati, a Northwestern Iranian language. It is spoken in Takestan, Qazvin province, Iran. Very little material exists. The searches found no primers, textbooks, dictionaries or audio corpora. The one substantial scholarly source is a 2024 phonetic description.

\section{Primary Educational Resources}
No primers, textbooks or school programmes were found in search results.

\section{Linguistic and Reference Materials}
\begin{itemize}
\item Taherkhani and Nelson, ``Southern Tati: Takestani dialect'', Journal of the International Phonetic Association 54(1), 2024 (DOI 10.1017/S0025100323000270). It is a phonetic description. \url{https://core-cms.cambridgecore.org/core/journals/journal-of-the-international-phonetic-association/article/southern-tati-takestani-dialect/9C2D1269604E1F8AAC18EAE6C7E92A5C}
\item The same article is indexed at the University of Illinois experts database. \url{https://experts.illinois.edu/en/publications/southern-tati-takestani-dialect/}
\item The same article is indexed at EBSCO. \url{https://www.ebsco.com/articles/social-sciences-and-humanities/b4dcfd73-26c7-5e6b-9da2-b72dbf9a741d/southern-tati-takestani-dialect}
\item Wikipedia entry on Tati (Iran), which lists Takestani among the varieties. \url{https://www.Wikipedia.org/wiki/Kabatei}
\end{itemize}

\section{Religious and Community Resources}
None were found.

\section{Digital and Audio Resources}
None were found.

\section{Curriculum Development Context}
Takestani is not known to be taught in schools. A curriculum would have to start from new documentation and community work. The 2024 JIPA article is a basis for orthography design. Its literature review cites earlier work by Yarshater (1969), Taherkhani (2009) and Rahmani and Rahmani (2021), which I did not locate online.
\end{document}